\section{公式插入}

公式插入示例如公式(\ref{E.example})所示。

薛定谔方程在一维量子问题中的应用：

（1）一维无限深势阱的势能分布

\begin{equation}
\displaystyle U(x)=\left\{
\begin{array}{l@{}}
0\;\;\quad(0<x<a) \\
U_0\quad (x\leq0 \;\mbox{或} \;x\geq a)\\ 
\end{array}\right.
\label{E.example}
\end{equation}

其本征能量与本征态为
\begin{equation}
\displaystyle U(x)=\left\{
\begin{array}{l@{}}
\displaystyle E_n=\frac{\hbar^2k^2}{2m}=\frac{\pi^2\hbar^2}{2ma^2}n^2 \;\;(n=1\mbox{，}2\mbox{，}3\cdots) \\
\displaystyle \psi _n(x)=\sqrt{\frac{2}{a}}\mathrm{sin}\frac{n\pi}{a}x \;\;\;(0\leq x\leq a) \\ 
\displaystyle \psi _n(x)=0\quad\quad\quad\quad\quad(x\leq0 \;\mbox{或} \;x\geq a) \\ 
\end{array}\right.
\end{equation}

（2）一维谐振子的势能函数$\displaystyle U(x)=\frac{1}{2}kx^2=\frac{1}{2}m\omega^2x^2$，谐振子的定态薛定谔方程为

\begin{equation}
\displaystyle -\frac{\hbar^2}{2m}\frac{\partial^2\psi}{\partial x^2}+\frac{1}{2}m\omega^2x^2\psi=E \psi
\end{equation}

其本征能量与本征态为
\begin{equation}
\displaystyle E_n=(n+\frac{1}{2})\hbar \omega \;\;(n=0\mbox{，}1\mbox{，}2\cdots)
\end{equation}
 
\begin{equation}
	\displaystyle \psi _n(\alpha x)=(\sqrt{\pi}2^nn!)^{-\frac{1}{2}}\mathrm{e}^{-\frac{\alpha^2x^2}{2}} H_n(\alpha x) \;\;(n=0\mbox{，}1\mbox{，}2\cdots)
\end{equation}

（3）一维方势垒和隧道效应：方势垒的势能为
\begin{equation}
	\displaystyle U(x)=\left\{
\begin{array}{l@{}}
U_0\quad(0<x<a) \\
0\;\;\quad( x\leq0 \;\mbox{或} \;x\geq a)\\ 
\end{array}\right.
\end{equation}

投射系数为$\displaystyle D=D_0\mathrm{e}^{-2ka}$，其中$\displaystyle k=\sqrt{\frac{2m(U_0-E)}{\hbar^2}}$。

我们还可以轻松打出一个矩阵
\begin{equation}
\bm{A}=\begin{bmatrix} %\bm为数学粗体 \hm为更粗的加重体数学负号
1&2&3&4\\
11&22&33&44\\
\end{bmatrix}
\times\begin{bmatrix}
22&24\\
32&34\\
42&44\\
52&54\\
\end{bmatrix}
\end{equation}
